% A new disjoint contribution to the selected M800 ideal fourth integral.
\subsection{A finite certificate for the unused fourth-history edge}

This computation implements the proved clipped historical-band method
on the actual mesh-$800$ outer history leaves, at $w=19/125$. The
chosen numerical fourth integral is the independently reconstructed
mesh-$800$ value, not an overlapping mesh-$400$ value. The old ordinary
compiler, ideal $C_3$, all-mask $C_4$, and common finite causal
approximation remain unchanged.

For each actual leaf $B=\prod_{j=1}^3[l_j,u_j]$, reconstruct exactly
the same constant bounds $L_0\le C_3\le U$ used by the old fourth
core and set $q_0=(L_0-\sum u_j)/3$. The old core has
$r<\min(l_3,q_0)$; its unique-lowest-bin ordering strip has
$r\le q_0$. In particular, when the new construction is nonempty,
$q_0<l_3$ and that leaf's old ordering strip is empty. Every new
interval starts at or above this actual $q_0$ and ends at or below
$\min_j l_j$. It is therefore disjoint from both old contributions.
No coarser ancestor's threshold is substituted.

The first minorant is the constant $L_0$. For the second, select a
branch $\min_j L_j(t)$ of the established symmetric third-cap
representation by a rational centre value. Choose one of its affine
planes $A$, and put
\[
 \kappa=\max\{0,\max_j\max_{t\in B}(A(t)-L_j(t))\},
 \qquad P=A-\kappa.
\]
The code asserts $P\le L_j$ throughout the full unsorted box for
every plane of that branch. Thus $P\le C_3$. Choosing the centre
only chooses a candidate; validity comes from these whole-box
inequalities. In particular, an active plane of a minimum is never
used directly as a minorant without this check and shift.

Let $\mathcal P$ be these one or two minorants and define
\[
 R=\min\left\{\min_j l_j,
       \max_{P\in\mathcal P}\frac13\max_{t\in B}(P(t)-S(t))\right\}.
\]
If $R\le q_0$ the leaf contributes nothing. Otherwise divide
$[q_0,R]$ into four equal rational intervals. On each one, denoted
$[r_0,r_1]$, test every mask containing the fourth coordinate.
For its selected old slots compute $b_-\le b(t)\le b_+$ and use
the candidate constant
\[
 L=\min\{(1+b_-+r_0)/2,57/64-3(b_++r_1)/2\}-2^{-40}.
\]
Retain it only if $L>U$ and every selected-remaining gate on the
entire unsorted box and entire interval is strictly below $L$.
All possible orders of old unselected slots are included on repeated
bins. Among eligible candidates retain one with largest $L$.
The endpoint buffer ensures that $L$ is strictly below the candidate
endpoint everywhere; other branch and history margins are treated
by the existing finite approximation.

For each $P\in\mathcal P$ the proved exact halfspace formula gives
\[
 V_P=\operatorname{Vol}\{(t,r)\in B\times[r_0,r_1]:
                          P(t)-S(t)-3r>0\}.
\]
Choose the larger volume, together with its actual minorant. These
alternative clipped sets can overlap and their volumes are not
added. The constant candidate $L$ is valid on the whole box and
interval independently of this choice, so the selected volume can
be combined with its common gap bound.

\paragraph{Outward gap calculation using the exact initial formula.}
Put $S_-=\sum l_j$ and
\[
 s_U=(U-S_--r_0)/r_0,\qquad s_L=(L-S_--r_0)/r_0.
\]
A positive clipped volume implies $s_U>2$. The common-corner
argument applies along the smaller rectangle leading from this
corner to each point of the clipped set, on which the constant-$U$
argument stays above two. Thus the actual fourth gap is at least
$\bar f(s_L)-\bar f(s_U)$ on that set, without requiring the
entire uncut band to have old argument above two.

Every actual leaf has total width at most $3/800$ and
$S_+\le3w$. Since $r_0\ge q_0\ge(\nu-S_+)/3$ and $L\le\Lambda$,
\[
 2<s_U<s_L\le
 \frac{3(\Lambda-3w+3/800)}{\nu-3w}-1
 =\frac{2791}{828}<4,
 \qquad \Lambda=153/256,\quad\nu=249/448.
\]
The inequality follows by maximizing
$3(\Lambda-S_++d)/(\nu-S_+)-1$ over $d\le3/800$ and
$S_+\le3w$; it is increasing in both variables. The exact initial
formula $\bar f(s)=\log(s-1)/s$ therefore applies to every counted
argument. Positive range-reduced atanh sums with their rational
tail bounds and outward integer divisions enclose both logarithms.
The counted gap floor is
$g=\lfloor Q(\bar f_{\rm lower}(s_L)-\bar f_{\rm upper}(s_U))\rfloor_+$,
with $Q=2^{112}$. In implementation both function bounds are already
in integer $Q$ units before this subtraction.

If $m_B$ is the product of factorials of repeated outer-bin
multiplicities, the recorded integer contribution is
\[
 \left\lfloor
 \frac{4gV_P}{m_B\,u_1u_2u_3r_1^2}
 \right\rfloor.
\]
Its value divided by $Q$ is a lower bound. The denominator uses
whole-band upper coordinates; the clipped volume is ordinary
Cartesian volume, not a product of logarithmic masses. Factorial
division is valid by the full-box minorant and gate checks and the
symmetry of the actual residual target. It is not an assertion that
the nonsymmetric clipped set has equal volume in all sorting sectors.

\paragraph{Cover, reproducibility, and addition rules.}
The pilot visits every $128$th previously positive parent and is an
overlapping diagnostic. The full run independently enumerates all
$294115$ actual history leaves of the chosen mesh and depth. It may
count a historical-edge band even where the earlier core happened
to have zero recorded lower gain. All old support is still excluded
by that leaf's $q_0$. An optional floating prefilter only omits leaves;
all counted minorants, levels, gates, volumes, logarithms and floors
are exact or outward-rounded.

Each positive witness row records its actual parent index and box,
band index, actual $q_0$, rational endpoints, mask and constant level,
minorant branch, shift and coefficients, exact clipped volume,
gap floor and final floor. Different bands within a leaf have
disjoint last-coordinate interiors; different leaves are disjoint
ordered sectors of the original history partition. Only one
minorant and one mask are counted per band. The full total may be
added to the selected mesh-$800$ fourth lower bound, while the pilot
must not also be added. The resulting new fourth total contributes
one quarter to its original $L_2$ occurrence.

The output is an author numerical certificate until separately
reconstructed. It proves an integral improvement under the existing
analytic inputs; it does not independently prove a distribution
theorem. Before applying the common finite approximation, trim the
clipped history wall and other equality boundaries by positive
margins and approximate the remaining compact sets by finitely
many disjoint boxes. The displayed strict buffers and any strict
global transfer margin absorb an arbitrarily small loss. All such
choices are made before $N$ tends to infinity.
